% Source: user supplied dfs1/dfs2 and path loops; ACM_Templates.pdf pp. 85-87.
% Supplement: https://cp-algorithms.com/graph/hld.html
\subsection{重链剖分：路径与子树区间}
\label{tree:hld}
\AlgorithmInfo
    {把树上路径分为少量连续区间，提供 LCA、子树区间和路径方向。}
    {预处理 $O(n)$；LCA、路径拆分 $O(\log n)$；存储 $O(n)$，路径结果 $O(\log n)$。}
    {非空连通树，指定根；这里只给剖分结构，不包含线段树等区间维护。}
\subsubsection*{原理}
第一次 DFS 计算父亲、深度、子树大小，把\textbf{子树点数最多}的儿子设为重儿子 $son_u$，其余为轻儿子。第二次 DFS 优先访问重儿子，继承链顶 $top$；轻儿子则另开新链。这样同一条重链在 DFS 序中连续，整棵子树也连续。

经过一条轻边向下时，子树大小至少减半，否则该儿子应是最大的儿子。因此根到任意点只经过 $O(\log n)$ 条轻边。处理 $u,v$ 的路径时，总是把链顶更深的一侧拆出，跳到该链顶的父亲，直到两点在同一链上，便只需 $O(\log n)$ 个区间。

\subsubsection*{使用教程}
树剖求 LCA 的跳链原理与选型说明见第~\ref{tree:hld-lca}~节。

\InlineCode{init(n)}、\InlineCode{add(u,v)}、\InlineCode{build(root)} 后，\InlineCode{fa,dep,sz,son,top,dfn,rk} 对应父亲、深度、大小、重儿子、链顶、序号与反向映射。\InlineCode{lca(u,v)} 可直接查询；\InlineCode{subtree(u)} 返回闭区间，区间中的序号通过 \InlineCode{rk} 还原为顶点。

\InlineCode{path(u,v)} 返回按原路径 $u\to v$ 顺序排列的区间 \InlineCode{Part{l,r,rev}}。每段均有 $l\le r$；\InlineCode{rev=true} 表示按 $r,r-1,\ldots,l$ 读取，否则按 $l,l+1,\ldots,r$ 读取。对求和等可交换操作，方向可以忽略；拼接、矩阵乘、颜色段合并等必须处理方向。
\begin{lstlisting}[style=mtycpp,numbers=none]
HLD g;
g.init(5);
g.add(1,2),g.add(1,3),g.add(2,4),g.add(2,5);
g.build(1);
int p=g.lca(4,3);
auto [l,r]=g.subtree(2);
auto parts=g.path(4,3);
vector<int> vertices;
for(auto [L,R,rev]:parts){
	if(rev) for(int i=R;i>=L;i--) vertices.push_back(g.rk[i]);
	else for(int i=L;i<=R;i++) vertices.push_back(g.rk[i]);
}
\end{lstlisting}
\InlineCode{path(u,v,true)} 切换为边模式，排除 LCA 对应位置；\InlineCode{subtree(u,true)} 返回子树内部边对应的区间，叶子会得到 $l>r$ 的空区间。具体边化点方法见第~\ref{tree:edge-to-vertex}~节。新增边或换根后重新构建，数组会重置。
\needspace{7\baselineskip}
\subsubsection*{模板代码}
\CodeFile{code/trees/hld.cpp}
